\section{Large biases}\label{sec:large-bias}

The last bias range admits a direct estimate. We retain the reduction
$0\le\rho\le1$. As throughout the proof, $H(t)=\sqrt{1-t^2}$,
expectation is uniform on the cube, and
\[
 m=\E f,\qquad d=T_\rho f,\qquad h=H(m),\qquad
 b=\E H(d),\qquad r=\rho^2=1-s^2.
\]
All terminating decimals in this section denote exact rational numbers.

\begin{proposition}\label{prop:large-bias}
If $f:\{-1,1\}^n\to\{-1,1\}$ has $0.74\le m<1$ and
$0.84\le s\le1$, then
\[
 h-b\le1-s.
\]
\end{proposition}

Write \(U=T_\rho\mathbf1_{\{f=-1\}}=(1-d)/2\), so that
\(H(d)=2\sqrt U\sqrt{1-U}\). Expanding \(\sqrt{1-U}\) will
reduce the desired lower bound on \(b\) to a lower bound on
\(\E U^{1/2}\) and upper bounds on the higher moments
\(\E U^{k+1/2}\), \(k\ge1\). These are supplied by reverse and
forward cube norm inequalities, respectively.

The forward inequality belongs to the classical hypercontractive
theory of Bonami and Beckner~\cite{Bonami70,Beckner75}; the reverse
inequality is due to Borell~\cite{Borell82}. We derive both from the
cube logarithmic Sobolev inequality using Gross's semigroup
method~\cite{Gross75}. The sharp reverse time constant is also given
by Mossel, Oleszkiewicz, and Sen~\cite[Equation~(1.2) and Theorem~1.10]{MOS13}
with logarithmic Sobolev constant $2$. The proof below fixes the
uniform-bit normalization and the limiting argument for inputs with zeros.

\subsection{The cube norm inequalities}

For a positive function $z$ on a finite probability space, write
\[
 \operatorname{Ent}(z)
 =\E(z\log z)-(\E z)\log(\E z).
\]
Recall the operators from Section~\ref{sec:preliminaries}: on the cube,
let $\E_i$ average only coordinate $i$, and set
$D_i=I-\E_i$ and $\mathcal N=\sum_iD_i$.
These self-adjoint projections give
\begin{equation}\label{lb:dirichlet}
 \E v\mathcal N w=\sum_i\E(D_iv)(D_iw).
\end{equation}
In particular, if $v_+,v_-$ and $w_+,w_-$ are the two values on a
coordinate edge, its conditional contribution is
\[
 \E_i[vD_iw]=\frac{(v_+-v_-)(w_+-w_-)}4.
\]

\begin{lemma}[Cube logarithmic Sobolev inequality]\label{lb:log-sobolev}
For every positive function $a$ on the uniform cube,
\begin{equation}\label{lb:entropy-bound}
 \operatorname{Ent}(a^2)\le2\E a\mathcal Na.
\end{equation}
\end{lemma}

\begin{proof}
First consider one bit. Both sides are homogeneous of degree two in
$a$. After dividing by $\E a^2$ and exchanging the two endpoints if
necessary, we may suppose that $a_\pm^2=1\pm x$, where $0\le x<1$.
The entropy then equals
\[
 J(x)=\frac{(1+x)\log(1+x)+(1-x)\log(1-x)}2,
\]
whereas the right side of \eqref{lb:entropy-bound} equals
\[
 \frac{(a_+-a_-)^2}{2}=1-\sqrt{1-x^2}.
\]
Both expressions vanish at zero. Their derivatives satisfy
\[
 J'(x)=\operatorname{arctanh}x\le\frac{x}{\sqrt{1-x^2}}.
\]
Indeed, the difference on the right minus the left vanishes at zero
and has derivative
$(1-x^2)^{-3/2}-(1-x^2)^{-1}\ge0$.
Integration proves the one-bit inequality, including its constant $2$.

For completeness, entropy is convex in its positive input. Its second
variation in the direction $w$ is
\[
 \E\frac{w^2}{z}-\frac{(\E w)^2}{\E z}\ge0
\]
by Cauchy--Schwarz. Let $\operatorname{Ent}_i$ denote entropy on bit $i$
with the other bits fixed. The exact decomposition
\[
 \operatorname{Ent}(z)
 =\E\operatorname{Ent}_1(z)+\operatorname{Ent}(\E_1z)
\]
and induction on the remaining coordinates imply
\begin{equation}\label{lb:tensorization}
 \operatorname{Ent}(z)\le\sum_i\E\operatorname{Ent}_i(z).
\end{equation}
To see the induction step explicitly, for $i\ne1$ convexity gives
$\operatorname{Ent}_i(\E_1z)\le\E_1\operatorname{Ent}_i(z)$:
conditional averaging in coordinate $1$ is a convex combination of
the two positive vectors indexed by coordinate $i$.
Apply the one-bit inequality to $z=a^2$ in every summand of
\eqref{lb:tensorization}, and then sum the edge energies using
\eqref{lb:dirichlet}.
\end{proof}

\begin{lemma}[Forward and reverse cube norm bounds]\label{lb:norm-bounds}
Let $u_0\ge0$ be a real function on the cube, $\tau\ge0$, and
$r=e^{-2\tau}$. For every $q>0$, $q\ne1$, put
$p_q=1+r(q-1)$. Then
\begin{equation}\label{lb:norm-moments}
 \E(T_{e^{-\tau}}u_0)^q
 \begin{cases}
 \displaystyle\le(\E u_0^{p_q})^{q/p_q},&q>1,\\[2pt]
 \displaystyle\ge(\E u_0^{p_q})^{q/p_q},&0<q<1.
 \end{cases}
\end{equation}
\end{lemma}

\begin{proof}
We begin with strictly positive $u_0$. The noise semigroup is
$T_{e^{-t}}=e^{-t\mathcal N}$, so $u(t)=T_{e^{-t}}u_0$ satisfies
$u'=-\mathcal Nu$. This identity follows, for example, because
$\mathcal N$ multiplies a character of degree $j$ by $j$, while noise
multiplies it by $e^{-tj}$.

For positive edge values $A>B$ and any $q>0$, $q\ne1$,
Cauchy--Schwarz in the interval $[B,A]$ gives
\begin{align*}
 \frac4{q^2}(A^{q/2}-B^{q/2})^2
 &=\left(\int_B^A t^{q/2-1}\,dt\right)^2\\
 &\le(A-B)\int_B^A t^{q-2}\,dt
 =\frac{(A-B)(A^{q-1}-B^{q-1})}{q-1}.
\end{align*}
The last quotient is positive also when $0<q<1$.
The inequality is symmetric in $A,B$ and holds by continuity when they
coincide. Summing it with the edge factor $1/4$ yields
\begin{equation}\label{lb:power-energy}
 \frac{\E u^{q-1}\mathcal Nu}{q-1}
 \ge\frac4{q^2}\E u^{q/2}\mathcal N(u^{q/2}).
\end{equation}

To reach a prescribed final exponent $q$, vary the exponent along
\[
 \theta(t)=1+(q-1)e^{2(t-\tau)},\qquad 0\le t\le\tau.
\]
This path is positive, stays on the same side of $1$ as $q$, and satisfies
$\theta'=2(\theta-1)$, $\theta(0)=p_q$, and $\theta(\tau)=q$.
For $M(t)=\E u(t)^{\theta(t)}$, direct differentiation gives
\[
 M'=\theta'\E(u^\theta\log u)
       -\theta\E u^{\theta-1}\mathcal Nu.
\]
Consequently, with $\|u\|_\theta=(\E u^\theta)^{1/\theta}$ also for
$0<\theta<1$,
\begin{align}
 \frac{d}{dt}\log\|u\|_\theta
 &=\frac{\theta'}{\theta^2}
      \frac{\operatorname{Ent}(u^\theta)}{M}
      -\frac{\E u^{\theta-1}\mathcal Nu}{M}\notag\\
 &=\frac{\theta-1}{M}
   \left[
    \frac2{\theta^2}\operatorname{Ent}(u^\theta)
    -\frac{\E u^{\theta-1}\mathcal Nu}{\theta-1}
   \right].\label{lb:norm-derivative}
\end{align}
Lemma~\ref{lb:log-sobolev}, applied to $u^{\theta/2}$, and
\eqref{lb:power-energy} show that the bracket is nonpositive.
Thus the derivative is nonpositive when $q>1$ and nonnegative when
$0<q<1$. Integrating, and raising the resulting norm inequality to
the positive power $q$, proves \eqref{lb:norm-moments}.
For general $u_0\ge0$, apply the result to $u_0+\varepsilon$ and let
$\varepsilon\downarrow0$. The cube is finite, and all exponents involved
are positive, so both sides converge continuously.
\end{proof}

\subsection{Reduction to a series of powers}

\begin{proof}[Proof of Proposition~\ref{prop:large-bias}]
We first use the norm inequalities to bound a positive series of moments.
We then bound its terms and tail uniformly on four intervals of $r$.
Finally, exact rational estimates make the resulting bound smaller than
$1-s$ on each interval.

The case $s=1$ has $\rho=0$ and $d=m$, hence $h-b=0$.
We may therefore assume $0.84\le s<1$, so $0<r\le0.2944$.
Write
\[
 p=\Pr(f=-1)=\frac{1-m}{2}\in(0,0.13],\qquad
 U=T_\rho\mathbf1_{\{f=-1\}}=\frac{1-d}{2}.
\]
The preceding approximation argument applies to this indicator input.
After taking its limit, $0\le U\le1$ and
$H(d)=2\sqrt U\sqrt{1-U}$.

For $0\le z\le1$, the binomial expansion has positive coefficients:
\begin{equation}\label{lb:binomial}
 1-\sqrt{1-z}=\sum_{k\ge1}c_kz^k,\qquad
 c_k=\frac{\binom{2k}{k}}{(2k-1)4^k},\qquad
 c_1=\frac12,\quad
 c_{k+1}=c_k\frac{k-1/2}{k+1}.
\end{equation}
The usual Taylor expansion gives this identity for $z<1$; positivity
and monotone convergence as $z\uparrow1$ give $\sum_{k\ge1}c_k=1$ and
the identity at $z=1$. It follows that
\begin{equation}\label{lb:series-identity}
 h-b
 =2(\sqrt p-\E\sqrt U)
  +2\sum_{k\ge1}c_k\bigl(\E U^{k+1/2}-p^{k+1/2}\bigr).
\end{equation}
The interchange with expectation is justified by the nonnegative
series for each of the two terms; both series are bounded by
$\sum_kc_k=1$.
Lemma~\ref{lb:norm-bounds} with $q=1/2$ bounds the first expectation
from below, and with $q_k=k+1/2$ bounds the other moments from above.
Thus
\begin{equation}\label{lb:series-bound}
 h-b\le
 2\bigl[p^{1/2}-p^{1/(2-r)}\bigr]
 +2\sum_{k\ge1}c_k
 \left[p^{q_k/(1+r(q_k-1))}-p^{q_k}\right].
\end{equation}

Every bracket on the right increases with $p$ on $[0,0.13]$.
Indeed, if $1/2\le a<b$ and $p>0$, then
\[
 \frac{d}{dp}(p^a-p^b)\ge0
 \quad\Longleftrightarrow\quad
 p\le(a/b)^{1/(b-a)}.
\]
Since $\log(b/a)\le(b-a)/a$, this threshold is at least
$e^{-1/a}\ge e^{-2}>0.13$; the last strict inequality follows also
from the logarithm bounds established in \eqref{lb:log-brackets} below.
The exponents in \eqref{lb:series-bound} have this ordering, since
$0<r<1$. We can therefore replace $p$ by
$p_*=13/100$ throughout its right side.

\subsection{Uniform bounds on the four noise intervals}

Put $L=\log(100/13)$ and $L_0=2041/1000$.
To replace powers of $p_*$ by rational expressions, define
\begin{equation}\label{lb:exponential-envelope}
 S_N(x)=\sum_{j=0}^N\frac{x^j}{j!},\qquad
 E(z)=\frac1{S_{40}((51/25)z)}\quad(z\ge0).
\end{equation}
We have the exact rational comparisons
\begin{equation}\label{lb:log-brackets}
 \begin{split}
 S_{15}\!\left(\frac{51}{25}\right)
 +\frac{(51/25)^{16}}{16!\bigl(1-(51/25)/17\bigr)}
 &<\frac{7691}{1000}<\frac{100}{13}\\
 &<\frac{7698}{1000}
 <S_{15}\!\left(\frac{2041}{1000}\right).
 \end{split}
\end{equation}
For $x=51/25$, the remainder after term $15$ in $e^x$ is at most
$x^{16}/[16!(1-x/17)]$: every ratio after the first omitted term is
at most $x/17<1$. The truncated sum at $2041/1000$ is less than its
full exponential. Therefore \eqref{lb:log-brackets} proves
\begin{equation}\label{lb:rational-envelopes}
 2.04<L<2.041=L_0,\qquad
 \sqrt{p_*}<0.361,\qquad p_*^z\le E(z)\quad(z\ge0).
\end{equation}
The middle inequality follows from $0.361^2>0.13$.
For the last one, $S_{40}(2.04z)\le e^{2.04z}\le e^{Lz}$;
no upper estimate for the omitted exponential tail is needed there.

Fix one of the four intervals
\begin{equation}\label{lb:noise-intervals}
 [\sigma_-,\sigma_+]
 \in\{[0.84,0.88],[0.88,0.92],[0.92,0.96],[0.96,1]\},
\end{equation}
and let
\[
 \ell=1-\sigma_+^2,\qquad R=1-\sigma_-^2.
\]
For $s$ in this interval and $s<1$, we have $0<r\le R$ and
$r\ge\ell$. Every expression with $\ell$ in its denominator below
is omitted when $\ell=0$.

First, put $\delta=r/[2(2-r)]$, so that $1/(2-r)=1/2+\delta$.
Convexity of $t\mapsto e^{-Lt}$ gives the trapezoid estimate
\[
 \frac{2[p_*^{1/2}-p_*^{1/(2-r)}]}{1-s}
 =\frac{2\sqrt{p_*}L}{1-s}\int_0^\delta e^{-Lt}\,dt
 \le\sqrt{p_*}L\frac{1+s}{1+s^2}
       \frac{1+p_*^\delta}{2}.
\]
Here we used $\delta/(1-s)=(1+s)/[2(1+s^2)]$.
The function $(1+s)/(1+s^2)$ decreases for $s\ge0.84$, since its
derivative has numerator $1-2s-s^2<0$.
Also $\delta$ increases with $r$, whence
$p_*^\delta\le p_*^{\ell/[2(2-\ell)]}$.
Using \eqref{lb:rational-envelopes}, define the resulting rational bound
\begin{equation}\label{lb:B}
 B=L_0(0.361)\frac{1+\sigma_-}{1+\sigma_-^2}
       \frac{1+E(\ell/[2(2-\ell)])}{2}.
\end{equation}

We next bound each remaining difference after division by $r$.
For $q=q_k=k+1/2$, put
\[
 z_q(t)=\frac q{1+t(q-1)},\qquad
 a=z_q(R),\qquad b_1=z_q(\ell),\qquad
 F_q(t)=p_*^{z_q(t)}-p_*^q.
\]
For every $q\ge3/2$ and $0\le t\le R\le184/625$, all the exponents
in use satisfy
\begin{equation}\label{lb:exponent-range}
 z_q(t)\ge\frac{3/2}{1+(184/625)/2}
 =\frac{625}{478}>1.3.
\end{equation}
To verify the first inequality, $z_q(t)$ decreases in $t$ and increases
in $q$, because $t<1$.
By $L>2.04$, the functions $z p_*^z$ and $z^2p_*^z$ strictly decrease
for $z\ge1.3$: their derivative signs are those of $1-Lz$ and
$2-Lz$, respectively.

There are three useful estimates for $F_q(r)/r$.
\begin{enumerate}
\item Differentiation gives
\[
 F_q'(t)=L(1-1/q)z_q(t)^2p_*^{z_q(t)}.
\]
Since $F_q(0)=0$ and $z_q(t)\ge a$ for $0\le t\le r$, integration
and the decreasing property just proved give
\[
 \frac{F_q(r)}r\le L_0(1-1/q)a^2E(a).
\]
\item If $\ell>0$, discarding the subtracted power and using
$r\ge\ell$ gives the direct estimate
\[
 \frac{F_q(r)}r\le\frac{p_*^{z_q(r)}}{\ell}
 \le\frac{E(a)}\ell.
\]
\item Integration in the exponent, followed by the trapezoid bound,
gives, with $z=z_q(r)$,
\begin{align*}
 \frac{F_q(r)}r
 &=\frac Lr\int_z^q p_*^t\,dt
 \le\frac{L(q-z)}{2r}(p_*^z+p_*^q)\\
 &=\frac{L(q-1)}2\bigl[zp_*^z+zp_*^q\bigr]
 \le\frac{L_0(q-1)}2\bigl[aE(a)+b_1E(q)\bigr].
\end{align*}
Here $q-z=r(q-1)z$, the first summand uses the decrease of $zp_*^z$
and $z\ge a$, and the second uses $z\le b_1$.
\end{enumerate}
It follows that, for $1\le k\le10$,
\begin{equation}\label{lb:Bk}
 \begin{split}
 \frac{F_{q_k}(r)}r\le B_k:=\min\biggl\{&
 L_0(1-1/q_k)a^2E(a),\;
 \frac{E(a)}\ell,\\
 &\frac{L_0(q_k-1)}2[aE(a)+b_1E(q_k)]
 \biggr\}.
 \end{split}
\end{equation}

For the tail, set
\begin{equation}\label{lb:Btail}
 a_{11}=\frac{23/2}{1+R(21/2)},\qquad
 B'=\min\left\{L_0a_{11}^2E(a_{11}),\;
                   \frac{E(a_{11})}\ell\right\}.
\end{equation}
For $k\ge11$, $z_{q_k}(R)\ge a_{11}$, again by its increase with $q$.
The first estimate for $F_q(r)/r$, with $1-1/q\le1$, is therefore
bounded by $L_0a_{11}^2E(a_{11})$, using the decrease of $z^2p_*^z$
before replacing the power by $E$.
The direct estimate is bounded by $E(a_{11})/\ell$ in the same way.
This proves $F_{q_k}(r)/r\le B'$ for every $k\ge11$.
In the interval ending at $1$, only the first of these tail estimates
is used; no division by zero occurs.

Finally, define
\begin{equation}\label{lb:finite-tail-sums}
 \begin{split}
 S&=\sum_{k=1}^{10}c_k B_k,\qquad
 T=\left(1-\sum_{k=1}^{10}c_k\right)B',\\
 C&=B+2(1+\sigma_+)(S+T).
 \end{split}
\end{equation}
Since $r/(1-s)=1+s\le1+\sigma_+$, Equations~\eqref{lb:series-bound},
\eqref{lb:B}, \eqref{lb:Bk}, and \eqref{lb:Btail} imply
\begin{equation}\label{lb:final-reduction}
 \frac{h-b}{1-s}\le C.
\end{equation}

\subsection{Exact rational evaluation}

Here are the coefficient data and arithmetic recipe for the final check.
The first ten coefficients from \eqref{lb:binomial}, and their remaining
mass, are exactly
\begin{align}
 (c_1,\ldots,c_{10})
 &=\left(\frac12,\frac18,\frac1{16},\frac5{128},\frac7{256},
 \frac{21}{1024},\frac{33}{2048},\frac{429}{32768},
 \frac{715}{65536},\frac{2431}{262144}\right),\notag\\
 1-\sum_{k=1}^{10}c_k&=\frac{46189}{262144}.
 \label{lb:coefficient-data}
\end{align}
For rational $x$, compute $S_N(x)$ entirely rationally by starting
$t_0=1$ and successively taking $t_j=t_{j-1}x/j$, then summing
$t_0+\cdots+t_N$. This evaluates the two order-$15$ comparisons in
\eqref{lb:log-brackets} and every order-$40$ expression in
\eqref{lb:exponential-envelope}. Substitution in
\eqref{lb:B}--\eqref{lb:finite-tail-sums}, using only rational
multiplication, addition, division, and comparison, gives
Table~\ref{lb:arithmetic-table}. A minimum in \eqref{lb:Bk} or
\eqref{lb:Btail} is found by comparing its rational candidates; the
candidate with denominator $\ell$ is excluded when $\ell=0$.

In the table, $\overline B=\lceil10^9 B\rceil$,
$\overline S=\lceil10^9 S\rceil$, and
$\overline T=\lceil10^9 T\rceil$ are \emph{integers}.
Thus the table also supplies the simpler rational upper bound
\begin{equation}\label{lb:rounded-total}
 C\le\frac{\overline B+
                2(1+\sigma_+)(\overline S+\overline T)}{10^9}.
\end{equation}
For an exact fraction $x=A/D$ with positive denominator, its required
ceiling is $\lceil10^9A/D\rceil$, evaluated by integer division.
The recipe therefore involves no floating-point operations, including
in the rounding step. The executable companion
\texttt{verification/large\_bias.py} checks these coefficient data,
the logarithm brackets, each table entry, and the final comparisons.

\begin{table}[htbp]
 \centering
 \small
 \setlength{\tabcolsep}{4pt}
 \begin{tabular}{@{}cccrrrc@{}}
  \toprule
  $[\sigma_-,\sigma_+]$ & $\ell$ & $R$
  & $\overline B$ & $\overline S$ & $\overline T$ & upper bound\\
  \midrule
  $[0.84,0.88]$ & $141/625$ & $184/625$ &
  $746521967$ & $60049597$ & $2525240$ & $0.982$\\
  $[0.88,0.92]$ & $96/625$ & $141/625$ &
  $748896408$ & $56019174$ & $1084482$ & $0.969$\\
  $[0.92,0.96]$ & $49/625$ & $96/625$ &
  $750555827$ & $51401364$ & $283275$ & $0.954$\\
  $[0.96,1]$ & $0$ & $49/625$ &
  $751524751$ & $46967364$ & $36923$ & $0.940$\\
  \bottomrule
 \end{tabular}
 \caption{Exact integer ceilings for the rational bounds. The expression
 in \eqref{lb:rounded-total} is strictly below the final column in every
 row.}
 \label{lb:arithmetic-table}
\end{table}

For explicit final arithmetic, the four right sides of
\eqref{lb:rounded-total} are, in order,
\begin{gather*}
 \frac{24545083853}{25000000000}<\frac{982}{1000},\qquad
 \frac{3025545147}{3125000000}<\frac{969}{1000},\\
 \frac{23828990297}{25000000000}<\frac{954}{1000},\qquad
 \frac{939541899}{1000000000}<\frac{940}{1000}.
\end{gather*}
Every displayed upper bound is strictly below $1$.
The four intervals cover $0.84\le s<1$, so
\eqref{lb:final-reduction} proves $h-b\le1-s$ throughout the claimed
range. The already treated endpoint $s=1$ completes the proof.
\end{proof}
